\begin{abstract}
Children's speech emotion recognition (SER) lags behind adult SER despite sharing model architectures and training paradigms. Whether this gap stems from distribution shift between children's speech corpora or from suboptimal architectural design has remained an open question, as prior work has not systematically isolated these factors. Here we present a controlled 192-experiment ablation study using WavLM Base across three corpora---C-BESD (children's acted), FAU Aibo (children's spontaneous), and IEMOCAP (adult acted)---to disentangle the effects of distribution shift from pooling strategy, layer fusion, adapter modules, data augmentation, and transfer learning. We find that distribution shift between corpora, not architecture, is the dominant performance bottleneck: the gap between in-domain children's acted speech (91.87\%) and spontaneous speech (67.05\%) persists across all architectural configurations. Zero-shot cross-corpus transfer collapses to 19--35\%, and transfer learning success depends primarily on the target domain rather than the source. Architectural factors---self-attention pooling provides modest gains (+13.4pp over mean pooling on acted speech), while adapter modules and layer fusion contribute negligibly. These findings indicate that future progress in children's SER requires addressing the acted-spontaneous distribution gap rather than pursuing incremental architectural refinements. All 192 experiment logs, model checkpoints, and evaluation code are publicly released.
\end{abstract}
